\documentclass{situationist-preprint}

\title{Moral Anachronism and Reasonable Knowledge}
\msid{ETH-23-a}
\msfamily{ETH-23}
\msdiscipline{Normative and Applied Ethics}
\mstype{Principal manuscript}
\mskeywords{moral anachronism; blame; moral ignorance; relativism of distance; culpable ignorance; moral progress; historical judgment}

\begin{document}
\maketitle

\begin{abstract}
May we judge past people by present standards? Two answers dominate. Presentism holds that we may, without qualification: slaveholders acted wrongly and are to be blamed as we would blame a slaveholder today. Historicism holds that we may not: past agents must be judged by the standards of their time, and judgments across a great distance of outlook lack application. We argue for a third answer. Past acts are to be evaluated by the correct standards, of which present standards, insofar as they record moral progress, are the best approximation we have; there is no other standard for wrongness. But blameworthiness depends on what the agent could reasonably have known, and what a past agent could reasonably have known is fixed by the evidence and arguments available to her. The allowance for epistemic distance therefore enters at the level of blame, not of wrongness. We show that this view avoids the objections to both rivals, answers Williams's relativism of distance and Rosen's scepticism about culpable ignorance, and implies that the allowance is often smaller than historicists suppose.
\end{abstract}

\section{The problem}

Aristotle defended natural slavery; medieval and early modern jurists defended the subjection of women; many respectable people in the nineteenth century defended the exclusion of the poor from political life. In each case, the practice is now judged wrong, and the question arises how to judge those who took part in it. The charge of \emph{moral anachronism} is that judging them by our standards is as misplaced as faulting a Roman engineer for not using steel.

Two answers are common. \emph{Presentism} holds that the standards we now accept are the ones by which all conduct is to be judged, and that past agents are as blameworthy as present agents who act as they did. \emph{Historicism} holds that each agent is to be judged by the standards of her own society, so that a slaveholder in a slaveholding society who complied with its norms did nothing for which she can be criticized. Williams gave historicism its most careful form in his \emph{relativism of distance}: when two outlooks are related only notionally, so that neither is a real option for those who hold the other, the vocabulary of appraisal, ``right'', ``wrong'', ``unjust'', loses its grip on the relation \citep[ch.~9]{Williams1985}. We cannot, Williams thought, intelligibly ask whether the life of a Samurai or of a medieval peasant was just.

Each answer has decisive costs. Presentism blames people for failing to know what they had no way of knowing, and so violates a condition on blame that it accepts elsewhere: no one is blamed for failing to foresee a consequence that no reasonable agent could have foreseen. Historicism makes the wrongness of slavery relative to the society that practised it, and so cannot say that the abolition of slavery was an improvement rather than a change. That is incompatible with the reality of moral progress, which MTE-19 argues is real and measurable as the expansion of the moral circle.

The recent literature has located the difficulty more precisely, in the theory of culpable ignorance. \citet{Calhoun1989} argued that in abnormal moral contexts, in which a wrong practice is widely accepted, reproach may be appropriate even when blame is not. \citet{MoodyAdams1994} argued that much moral ignorance is affected: agents who could have known that their practices were wrong chose not to ask. \citet{Rosen2003} argued in the opposite direction that ignorance is culpable only if it results from an earlier culpable act, and that, since such acts are rarer than we suppose, many past wrongdoers are blameless. \citet{Fricker2010} defended a version of Williams's view restricted to blame. We argue for a view that agrees with Fricker about the locus of the relativity and with Moody-Adams about its extent.

\section{Framework}

\begin{definition}[Two appraisals]\label{d:appraisals}
The \emph{appraisal of an act} is the judgment whether it was right or wrong. The \emph{appraisal of an agent} is the judgment whether, and to what degree, she was blameworthy for the act.
\end{definition}

\begin{definition}[Reasonable knowledge]\label{d:reasonable}
An agent \emph{could reasonably have known} that an act was wrong if the evidence and arguments available to her, including those she was in a position to acquire by inquiry that she had reason to undertake, justified her in believing that it was wrong.
\end{definition}

The notion of justification in Definition~\ref{d:reasonable} is internalist. On the view defended in EPI-06-a, justification answers the question whether a subject has believed as she should, given what she has to go on, and is shared by internal duplicates; knowledge, by contrast, is externalist. This is the right notion for blame. Blame concerns whether the agent has acted as she should, given what she had to go on, and a notion of reasonable knowledge that depended on facts inaccessible to her would make her blameworthy for features of her situation that she had no means of registering.

\begin{definition}[Epistemic distance]\label{d:distance}
The \emph{epistemic distance} between a past agent and a moral truth is the degree to which the evidence and arguments available to her failed to justify belief in it.
\end{definition}

\begin{principle}[One standard of wrongness]\label{p:standard}
Past acts are right or wrong by the same standard as present acts. Present standards are to be applied to the past to the extent that they are better approximations to that standard than past standards were.
\end{principle}

\begin{principle}[Culpability]\label{p:culp}
An agent is blameworthy for a wrong act only to the degree that she could reasonably have known that it was wrong, or that her ignorance of its wrongness was itself something she could reasonably have known to be culpable.
\end{principle}

Principle~\ref{p:standard} is supported by the reality of moral progress. If moral progress is real, as MTE-19 argues, then later standards can be better than earlier ones, and ``better'' requires a standard against which both are measured. The same standard applies to the acts of those who held the earlier standards. Principle~\ref{p:culp} is a version of the epistemic condition on responsibility that goes back to Aristotle's treatment of ignorance as a ground of the involuntary \citep[III.1]{AristotleNE}. Its second clause allows for culpable ignorance without requiring, as Rosen does, an earlier act of akrasia.

\section{The argument}

\begin{theorem}\label{t:main}
We may judge past people by present standards in the appraisal of acts, without allowance for epistemic distance. In the appraisal of agents, we may judge them by present standards only with allowance for epistemic distance: a past agent is blameworthy for a wrong act only to the degree that she could reasonably have known it to be wrong.
\end{theorem}

\begin{proof}
For acts: by Principle~\ref{p:standard}, a past act is wrong if it is wrong by the correct standard, and present standards, where they record progress, approximate that standard better than the standards of the agent's time. Wrongness is a property of the act, and it does not depend on what the agent knew. So the appraisal of a past act by present standards is legitimate, and the agent's epistemic situation is irrelevant to it.

For agents: by Principle~\ref{p:culp}, blameworthiness requires that the agent could reasonably have known that the act was wrong. By Definition~\ref{d:reasonable}, this depends on what her available evidence and arguments justified, which varies with her historical situation. By Definition~\ref{d:distance}, the greater the epistemic distance, the less the agent could reasonably have known, and so the less blameworthy she is.
\end{proof}

\begin{corollary}\label{c:neither}
Both presentism and historicism are false. Presentism is false of the appraisal of agents; historicism is false of the appraisal of acts.
\end{corollary}

The theorem identifies where the allowance for distance belongs, but it does not say how large it is. That is a historical question, and the answer is often that it is small. Moral progress, on the account in MTE-19, consists largely in the expansion of the moral circle to include those previously excluded, and the arguments for each expansion have usually been available, and often stated, well before the expansion occurred. Aristotle reports that some in his time held that mastery over slaves is contrary to nature and rests only on force \citep[I.3, I.6]{AristotlePol}. The argument against the justice of slavery was therefore available to him, and so was evidence, in the slaves he lived among, of the very capacities that he denied them. His epistemic distance from the wrongness of slavery was not great, and the allowance it earns him is correspondingly limited; ETH-23-b examines the case in detail. In the debates over the conquest of the Americas, the case against the enslavement of the indigenous peoples was argued publicly in Spain in the middle of the sixteenth century. And abolitionist arguments were in wide circulation in the societies that practised chattel slavery in the eighteenth and nineteenth centuries. In such cases the evidence available to a reflective agent justified the belief that the practice was wrong, and the allowance for distance does not exculpate.

The allowance does real work in other cases: those in which the relevant evidence was not available. An agent who acts on a scientific belief that was justified in her time, and whose act is wrong because the belief was false, is excused to the degree that the belief was justified. The same holds for moral beliefs whose correction required conceptual resources, such as the concept of a structural injustice, that were not available to her. The more the progress we credit depended on new arguments, the greater the allowance; the more it depended on the application to new cases of arguments already accepted, the smaller.

\section{Objections and replies}

\begin{objection}[The relativism of distance]
Williams held that across a notional confrontation the vocabulary of appraisal lacks application. To call an ancient practice unjust is not false but idle, since no one could take up or reject the ancient outlook as a real option.
\end{objection}

\begin{reply}
Williams's argument concerns the point of appraisal, and it is strongest for the appraisal of agents, where we agree that the agent's distance matters. For acts, the reality of moral progress refutes it. If the abolition of slavery was progress, then slavery was worse, by a standard that applies to both, than its abolition, and the appraisal of the ancient practice as unjust is the same judgment that makes abolition progress. \citet{Fricker2010} drew a similar distinction, retaining Williams's relativity for blame. Our Principle~\ref{p:culp} captures what is correct in this: the appropriate degree of blame depends on what was available to the agent. What it does not capture is a relativity of wrongness, which we deny.
\end{reply}

\begin{objection}[Rosen's regress]
On \citeauthor{Rosen2003}'s view, ignorance is culpable only if it results from an earlier culpable act or omission, done in knowledge of its wrongness. Few past agents committed such acts with respect to their moral beliefs. So few were blameworthy, and the allowance for distance is nearly total \citep{Rosen2003}.
\end{objection}

\begin{reply}
The regress arises because Rosen ties culpability to knowledge at the moment of acting. Principle~\ref{p:culp} ties it to reasonable knowledge, defined in terms of internalist justification. If the evidence available to an agent justified the belief that her act was wrong, then her failure to believe it was a failure to believe as she should, and that failure is itself a ground of blame; no prior akratic act is needed. The second clause of Principle~\ref{p:culp} does not require tracing to an earlier knowing wrong, only that the agent could reasonably have known that she should inquire. Rosen's view correctly identifies that blame requires a failing on the agent's part; it misidentifies the failing as necessarily akratic.
\end{reply}

\begin{objection}[Our standards may be wrong too]
Present standards are themselves the product of a contingent history, and future people may judge them as we judge Aristotle's. Applying them to the past is presumptuous.
\end{objection}

\begin{reply}
The objection supports Principle~\ref{p:standard} rather than undermining it. That future people may judge us correctly presupposes a standard that applies to us whether or not we accept it. Present standards are to be applied only insofar as they record progress, and MTE-19 holds that progress, though real, is neither steady nor irreversible; some past agents saw some things more clearly than we do. The view therefore recommends humility about the content of our standards, not about the existence of a standard. It also recommends applying Principle~\ref{p:culp} to ourselves: where the evidence available to us justifies the belief that a present practice is wrong, we will have no excuse.
\end{reply}

\begin{objection}[Moral luck]
The view makes blameworthiness depend on when one was born, which is a matter of luck. Two agents who acted alike, with the same character, would differ in blameworthiness because one lived in a society where the arguments were available and the other did not.
\end{objection}

\begin{reply}
They would, and this is a case of circumstantial moral luck. ETH-13 argues that moral luck is real and does not undermine responsibility. The present case is an instance: the circumstance is epistemic, and it affects what the agent could reasonably have known, and through that, how far she is to blame. It does not affect the wrongness of the act, which is the same for both.
\end{reply}

\section{Conclusion}

We may judge past people by present standards, but the two appraisals must be kept apart. Acts are right or wrong by a single standard, and present standards, where they reflect moral progress, are the best approximation to it we have. Agents are blameworthy only for what they could reasonably have known, and what a past agent could reasonably have known was fixed by the evidence and arguments available in her time. The allowance for epistemic distance is real, but in the central cases of moral progress it is smaller than historicists suppose, since the arguments for including the excluded were usually already in circulation.

\begin{thebibliography}{99}\small
\bibitem[Aristotle(NE)]{AristotleNE} Aristotle. \emph{Nicomachean Ethics}. Cited by book and chapter.
\bibitem[Aristotle(Pol.)]{AristotlePol} Aristotle. \emph{Politics}. Cited by book and chapter.
\bibitem[Calhoun(1989)]{Calhoun1989} Calhoun, C. (1989). Responsibility and reproach. \emph{Ethics}, 99(2).
\bibitem[Fricker(2010)]{Fricker2010} Fricker, M. (2010). The relativism of blame and Williams's relativism of distance. \emph{Proceedings of the Aristotelian Society, Supplementary Volume}, 84.
\bibitem[Moody-Adams(1994)]{MoodyAdams1994} Moody-Adams, M.~M. (1994). Culture, responsibility, and affected ignorance. \emph{Ethics}, 104(2).
\bibitem[Rosen(2003)]{Rosen2003} Rosen, G. (2003). Culpability and ignorance. \emph{Proceedings of the Aristotelian Society}, 103.
\bibitem[Williams(1985)]{Williams1985} Williams, B. (1985). \emph{Ethics and the Limits of Philosophy}. Fontana.
\end{thebibliography}

\end{document}
